% ID: SEL-D0003-OB-E04
% ESTADO: cotejada. Fuente y fragmento cotejados de forma independiente; salida compilada revisada.
% Fuente: Selectividad/sel_2019_matematicas_aplicadas/Suplente Junio/Exámenes Matemáticas Aplicadas CCSS - A y B.pdf
% SHA-256: a59474e6895bd41034da53b7060969cbe8cbf7f8272f39fbf16821c95f06c84b
% Página: 2; opción B; ejercicio 4.
% Transcripción de primera pasada basada en la imagen de la página original.
% Se conserva el espacio de millares, los ordinales y los puntos decimales.
% Sin solución ni modificación de datos o condiciones.

La Consejería de Educación elige una muestra de 5\,000 estudiantes de 1\textsuperscript{o} de Bachillerato de Ciencias Sociales y los encuesta para conocer la opinión que tienen sobre la elección de cierta materia entre las optativas para cursar 2\textsuperscript{o} de Bachillerato. El resultado de la encuesta revela que 2\,250 estudiantes piensan elegir dicha materia optativa.
\begin{enumerate}[label=\alph*),leftmargin=2.4em]
  \item \textbf{(1.5 puntos)} Halle un intervalo de confianza al 97.5\,\% para estimar la proporción de estudiantes que piensan elegir esa materia optativa.
  \item \textbf{(1 punto)} Si en otra muestra la proporción de estudiantes que piensa elegir esa materia es de 0.5 y el error cometido en la estimación ha sido inferior a 0.03 con un nivel de confianza del 92.5\,\%, calcule el tamaño muestral mínimo de esa muestra.
\end{enumerate}
